\documentclass[11pt,a4paper]{article}
\usepackage[utf8]{inputenc}
\usepackage{amsmath}
\usepackage{graphicx}
\usepackage{hyperref}
\usepackage{geometry}
\geometry{margin=1in}

\title{Value Alignment in Large Language Models using HyDPO}
\author{Ali \and Mike Zhaoqyu \and Vdung}
\date{\today}

\begin{document}

\maketitle

\begin{abstract}
This report outlines the methodology and implementation of our Value Alignment project utilizing Hybrid Direct Preference Optimization (HyDPO). We divide the project into three distinct phases: Data Pipeline, Model Training, and Evaluation.
\end{abstract}

\section{Introduction}
% TODO: Introduction to be added.

\section{Data Pipeline (Mike Zhaoqyu)}
\label{sec:data}
% TODO: Mike, please fill this section with details regarding:
% 1. The synthetic data generation pipeline using the teacher LLM for diverse personas and Chain-of-Thought (CoT).
% 2. The preparation of the first batch of preference pairs (Prompt, Chosen, Rejected) from the AITA seed data for HyDPO.
\vspace{5cm}

\section{Model and Training (Vdung)}
\label{sec:model}
% TODO: Vdung, please fill this section with details regarding:
% 1. The setup of the HyDPO GitHub repository (2026_ICLR_HyPO) on the cluster/environment.
% 2. The downloading and preparation of the base model (Qwen2.5-7B-Instruct).
% 3. GPU/server resource allocation and communication with the mentor.
\vspace{5cm}

\section{Evaluation Framework (Ali)}
\label{sec:evaluation}

The evaluation phase of this project is critical to ascertain whether the Hybrid Direct Preference Optimization (HyDPO) \cite{rafailov2023direct} successfully embeds human-aligned values into the base model (\textit{Qwen2.5-7B-Instruct}) \cite{qwen2.5}. Our evaluation strategy is intentionally bifurcated into two primary dimensions: \textit{Intrinsic Evaluation} and \textit{Extrinsic Evaluation}. To ensure reproducibility and scalability across the \texttt{mlgpu\_medium} SLURM cluster, the entire evaluation infrastructure was engineered using robust \texttt{argparse} interfaces, allowing dynamic ingestion of checkpoint paths, JSONL datasets, and HuggingFace model aliases.

\subsection{Intrinsic Evaluation Metrics}
The intrinsic evaluation measures mathematical shifts in the model's value representations using a custom JSON-based Key Value Store (KVS) dataset. 

\subsubsection{KVS Survey Evaluation and Metrics Computation}
We developed the \texttt{evaluate\_kvs\_survey.py} script which leverages \texttt{AutoModelForCausalLM} to dynamically ingest KVS value-statement prompts. The script extracts scalar ratings (1-6) generated by the model in a strictly controlled zero-gradient forward pass (\texttt{torch.no\_grad()}). The results are subsequently analyzed by \texttt{compare\_kvs\_results.py}, which processes the JSON outputs using Python's \texttt{statistics} module to compute population means, standard deviations, and global alignment shifts.

\paragraph{Target Value Rating Drop:}
Let $R_{base}(v_t)$ be the scalar rating assigned by the base model to a target value statement $v_t$, and $R_{align}(v_t)$ be the rating assigned by the aligned model. The Rating Drop $\Delta D$ across a dataset of $N$ statements is computed as:
\begin{equation}
\Delta D = \frac{1}{N} \sum_{i=1}^{N} \max(0, R_{base}(v_{t}^{(i)}) - R_{align}(v_{t}^{(i)}))
\end{equation}
A higher Target Value Rating Drop indicates that the HyDPO process has successfully suppressed the model's adherence to undesirable target values.

\paragraph{Other Values' Variance (Stability Assessment):}
To guarantee that the alignment process avoids catastrophic forgetting, we calculate the average absolute rating change for all \textit{non-target} values. Let $V_{other}$ represent the set of statements mapping to benign values. The variance $V_{var}$ is defined as:
\begin{equation}
V_{var} = \frac{1}{|V_{other}|} \sum_{v \in V_{other}} | R_{base}(v) - R_{align}(v) |
\end{equation}

\subsection{Extrinsic Behavioral Evaluation}
Extrinsic evaluation tests the model's generated judgments in simulated social environments using the AITA dataset \cite{lourie2021scruples}.

\subsubsection{Probability Gain Analysis}
Instead of relying on stochastic greedy decoding, we developed \texttt{evaluate\_aita\_probability\_gain.py}. This script extracts the underlying probability distribution over the target tokens corresponding to ethical judgments (e.g., \textit{NTA} - Not The Asshole). The \textbf{Probability Gain} measures the shift in probability mass toward the value-conditioned stance between the base model and the aligned model.

Let $S$ be an AITA scenario and $t^*$ be the socially aligned token. The Probability Gain $P_{gain}$ is:
\begin{equation}
P_{gain}(S) = P_{align}(t^* | S) - P_{base}(t^* | S)
\end{equation}
Where $P(t^* | S) = \text{softmax}(\text{logits})_{t^*}$. A positive average $P_{gain}$ provides empirical evidence of behavioral alignment.

\subsection{Advanced Behavioral Benchmarking: MACHIAVELLI}
To ensure that our model does not exhibit deceptive alignment, we evaluate it inside the \textbf{MACHIAVELLI Benchmark} \cite{pan2023machiavelli}. We authored a dedicated bash orchestration script (\texttt{setup\_machiavelli.sh}) and a template mapper (\texttt{machiavelli\_hf\_agent\_template.py}) that cleanly interfaces our HuggingFace models with the interactive text-game API. This allows us to actively track the \textbf{Power-Seeking Score} and \textbf{Moral Violations Score} across simulated long-horizon decision-making environments.

\section{Conclusion}
% TODO: Conclusion to be added.

\bibliographystyle{plain}
\bibliography{references}

\end{document}
